\chapter{Modelo neuronal: DeBERTa-v3 \emph{fine-tuned}}\label{cap:deberta}

\section{Arquitectura}

DeBERTa-v3-base \citep{he2023debertav3} es un modelo \emph{transformer}
de codificador con 12 capas, 768 dimensiones por capa, 12 cabezales de
atención y aproximadamente 184\,M parámetros. Las dos contribuciones
arquitectónicas centrales son:

\begin{enumerate}[leftmargin=*,itemsep=0.2em]
  \item \textbf{\emph{Disentangled attention}}: la matriz de atención
        se descompone en cuatro términos correspondientes a las
        interacciones contenido--contenido, contenido--posición,
        posición--contenido y posición--posición. Cada componente
        emplea proyecciones independientes:
        \[
          A_{ij} = Q^c_i K^{c\top}_j + Q^c_i K^{p\top}_{j|i} +
                   K^c_j Q^{p\top}_{i|j}.
        \]
  \item \textbf{Preentrenamiento RTD} (\emph{Replaced Token
        Detection}): un generador propone sustituciones léxicas y un
        discriminador (el modelo principal) clasifica cada \emph{token}
        como original o reemplazado. Esta señal binaria por
        \emph{token} resulta más eficiente que el clásico
        \emph{Masked Language Modeling}.
\end{enumerate}

\section{Adaptación a NLI sobre evidencias}

La entrada al modelo es la cadena tokenizada
\[
  [\text{CLS}]\, e \,[\text{SEP}]\, c \,[\text{SEP}],
\]
truncada a un máximo de 256 \emph{tokens} (la mayoría de los pares
encajan en este límite; los casos restantes se truncan eliminando los
\emph{tokens} finales de la evidencia). Sobre la representación del
\emph{token} \texttt{[CLS]} se aplica una capa lineal con $K = 3$
salidas, seguida de \emph{softmax}.

\section{Procedimiento de \emph{fine-tuning}}

\paragraph{Optimizador y \emph{schedule}.} AdamW
\citep{loshchilov2019decoupled} con $\beta_1 = 0.9$, $\beta_2 = 0.999$,
$\epsilon = 10^{-8}$ y \emph{weight decay} $= 0.01$ (excluyendo
\emph{bias} y parámetros de \emph{layer norm}). Tasa de aprendizaje
máxima $2 \cdot 10^{-5}$ con calentamiento lineal del 6\,\% de los
pasos totales y decaimiento lineal posterior.

\paragraph{Otros hiperparámetros.} Tamaño de lote efectivo 32
(\emph{gradient accumulation} cuando la memoria de GPU lo requiera),
3 épocas, precisión mixta \texttt{fp16}, \emph{label smoothing}
$\epsilon = 0.05$.

\paragraph{Reproducibilidad.} Las semillas (\texttt{42, 1337, 2025}) se
fijan en \texttt{torch}, \texttt{numpy} y \texttt{random}, y se activa
\texttt{torch.use\_deterministic\_algorithms(True)} en la medida en que
las operaciones de cuDNN lo permiten.

\section{Implementación}

Pipeline en \path{research-stats/research_stats/transformer/finetune.py}
empleando la API \texttt{Trainer} de \texttt{transformers}. La
exportación a ONNX para servir el modelo en producción se realiza con
\texttt{optimum} en \path{research_stats/export/to_onnx.py}.

\section{Resultados}

% TODO(resultados): completar tras entrenamiento real.
\begin{table}[h]
  \centering
  \begin{tabular}{lcccc}
    \toprule
    Modelo & Acc. & F1 macro & NLL & ECE \\
    \midrule
    TF-IDF + LR & --- & --- & --- & --- \\
    DeBERTa-v3 \emph{fine-tuned} & --- & --- & --- & --- \\
    \bottomrule
  \end{tabular}
  \caption{Comparativa baseline vs.\ modelo neuronal sobre la partición
  \emph{dev} de FEVER (media $\pm$ desv.\ típica, 3 semillas).}
  \label{tab:deberta-vs-baseline}
\end{table}

\section{Discusión preliminar}

Se discutirán los \emph{trade-offs} entre coste computacional
(entrenamiento e inferencia) y ganancia en calidad. La cuantificación
formal de la diferencia se aborda en el capítulo~\ref{cap:tests}.
